% Generated by make_glossary.py from the term list; do not edit by hand.
\begingroup\small
\par\noindent\textbf{HDL.} Hardware description language: the language, such as Verilog, in which a designer writes a circuit~\cite{jamiesonthesis}. % ref:hdl
\par\noindent\textbf{RTL.} Register-transfer level: the HDL text that says what the registers hold and how they change~\cite{jamiesonthesis}. % ref:rtl
\par\noindent\textbf{gate.} A basic logic element (AND, OR, NOT); the elaborated design is first written with such primitive gates~\cite{jamiesonthesis}. % ref:gate
\par\noindent\textbf{netlist.} The list of the basic blocks of a design and the nets (wires) that join them; every net is driven by exactly one block~\cite{luuthesis}. % ref:netlist
\par\noindent\textbf{flip-flop.} A one-bit storage element that keeps its value between clock edges. % def:ff
\par\noindent\textbf{LUT.} Look-up table: a small memory that implements any logic function of $K$ inputs ($K=6$ here). In most FPGAs a LUT and a flip-flop form one basic logic element~\cite{jamiesonthesis}~\cite{luuthesis}. % def:lut
\par\noindent\textbf{technology mapping.} The CAD step that covers the logic of a netlist with the cells the target offers: gates for a gate library, LUTs and flip-flops for an FPGA~\cite{jamiesonthesis}. % def:techmap
\par\noindent\textbf{provenance.} The link from each netlist object back to the HDL that produced it. % ref:prov
\par\noindent\textbf{Odin III.} The provenance-tracked synthesis front end under Loki; it keeps that link in one intermediate representation. % ref:odin
\par\noindent\textbf{answer key / ground truth.} What the synthesis tool recorded about which netlist objects form which HDL construct; the target that recovery is scored against. % ref:key
\par\noindent\textbf{Gate-Breaker.} An earlier study in which a language model rewrites gate netlists into RTL; Loki's method baseline. % ref:gb
\par\noindent\textbf{hypothesis.} A claim that an experiment can confirm or fail to confirm. In L1 a structure hypothesis is the claim that a set of netlist cells is a structure of some kind (register, adder, mux, ...). % ref:hyp
\par\noindent\textbf{L0 / L1 / L2.} Loki's three layers: L0 labelled datasets, L1 a Bayesian flow over structure hypotheses, L2 graph analysis. % ref:layers
\par\noindent\textbf{quartile.} One quarter of the benchmark designs ordered by Gate-Breaker's complexity score (Q1 simplest, Q4 most complex). % def:quart
\par\noindent\textbf{AST / CFG similarity.} How alike two RTL texts are by subtree overlap of the syntax trees (AST) and by a kernel on the control-flow graphs (CFG); Gate-Breaker's measures of a recovered text against its reference. % def:astcfg
\par\noindent\textbf{micro designs.} The 68 small synthesizable RTL designs (counters, registers, arithmetic, ALUs, muxes, comparators, small state machines) used as the first benchmark. % def:micro
\par\noindent\textbf{VerilogEval references.} The 156 reference RTL designs of the VerilogEval spec-to-RTL set; the second benchmark. % def:veval
\par\noindent\textbf{floor.} The abstraction level and technology mapping at which a design is cut and written as the netlist to reconstruct (f1-k6, f1-earch, f2, f3). % def:floor
\par\noindent\textbf{renamed-instance protocol.} Rename every identifier through one shared permutation, run the pipeline, map the tokens back: equal results show that the pipeline does not use names. % def:renamed
\par\noindent\textbf{detector.} A plug-in that proposes structure hypotheses for a netlist and reports features of them, never a likelihood. % def:det
\par\noindent\textbf{calibration.} Fitting priors and likelihood ratios on a training split so that predicted probabilities match observed rates. % def:cal
\par\noindent\textbf{ECE.} Expected calibration error: the average gap between predicted probability and observed match rate over ten bins. % def:ece
\par\noindent\textbf{prover.} An engine (cone check, truth table, transfer) that proves a structure hypothesis equal to the netlist; a proof makes it near certain. % def:prover
\par\noindent\textbf{count / match (scoring).} A hypothesis is counted when it is committed or has probability at least 0.5; it matches an answer-key instance by an overlap rule (Section~\ref{sec:eval}). % def:count
\par\noindent\textbf{precision / recall / F1.} Precision is the share of counted hypotheses that match; recall is the share of answer-key instances matched; F1 is their harmonic mean. % def:prf
\par\noindent\textbf{macro-F1 / micro-F1.} F1 averaged over the kinds of structure (macro) or pooled over all instances (micro). % def:macro
\par\noindent\textbf{Louvain / modularity / NMI.} Louvain is a community-detection method; modularity scores how dense the clusters are; NMI scores the agreement of a clustering with the known module hierarchy. % def:louv
\par\noindent\textbf{GNN and 1-WL baseline.} A graph neural network (GNN) labels each netlist node from its neighbours; the 1-WL baseline is colour refinement, whose distinguishing power a message-passing GNN cannot exceed. % def:gnn
\par\noindent\textbf{word.} A group of cells that act as one multi-bit signal (a register, a datapath operator, a bus). % def:word
\par\noindent\textbf{operator type-F1.} F1 of the recovered word-level operator types (adder, mux, comparator, ...) against the answer key's operators. % def:opf1
\par\noindent\textbf{fragment.} A wrong-looking candidate that is in fact part of a real construct. % def:frag
\par\noindent\textbf{word-first detector.} A detector that first groups flip-flops into register words by shared clock, enable and reset, then builds the rest around them. % def:wf
\par\noindent\textbf{DANA (through HAL).} A published method, run through the HAL framework, that groups flip-flops into register words from their shared control nets and the dataflow between them. % def:dana
\par\noindent\textbf{McNemar / Wilcoxon test.} An exact paired test on discordant pairs; a paired rank test. % def:tests
\par\noindent\textbf{synthetic corpus.} A seeded generator's 1200 labelled idiom designs (twelve families) with train, iid and ood splits. % def:corpus
\par\noindent\textbf{signature dictionary.} A table from a cell's cone class and WL colour to the construct it usually is, mined from the synthetic corpus. % def:dict
\par\noindent\textbf{D14 provenance map.} Odin III's version-2 name map: per object, the source-statement units it blends, and labelled join records. % def:d14
\par\noindent\textbf{unit, separable, lost.} A unit is the statement-level anchor of a construct; it is separable while some object has exactly that unit set, and lost at the first later snapshot where none has. % def:unit
\par\noindent\textbf{random-set null / persistence null.} The random-set null draws random unit sets of the real sizes at each pass; the persistence null starts from the objects of the previous snapshot and merges at random until the real object count is reached. % def:nullm
\par\noindent\textbf{LLM pilot arms A, B, C.} Arm A: gate netlist; arm B: netlist plus Loki hints; arm C: interface line only (memorisation control). % def:arms
\endgroup
